\section{Exact-degeneracy mechanisms and model families}
\label{sec:models}

The models are grouped by the algebra that fixes the protected rank, rather
than by spatial dimension or by the statistic that is eventually measured.
This organization is necessary because the same dimension $D$ can arise from
a common kernel, a cohomology group, or dependent stabilizer constraints, and
these mechanisms impose different restrictions on the response matrices.  In
every numerical case we identify the kernel first, verify a positive gap to
its complement, and only then construct tangent responses.  The ranks and gap
audits used in the comparison are imported from the generated evidence table,
not transcribed into the prose.

\subsection{Constraint kernels and clustered Hall parents}

A frustration-free parent can be written in terms of a constraint map $C$
as
\begin{equation}
 H_{\mathrm{par}}=C^\dagger C
 =\sum_\alpha C_\alpha^\dagger C_\alpha,
 \qquad
 \ker H_{\mathrm{par}}=\bigcap_\alpha\ker C_\alpha.
 \label{eq:factor-parent}
\end{equation}
Positivity pins every state in this common kernel to exactly zero energy.  It
does not fix the orientation of the kernel when the constraints change.  If
$C(\boldsymbol\lambda)$ has constant nullity and its smallest nonzero singular
value remains positive, Eq.~\eqref{eq:factor-parent} supplies precisely the
constant-rank, gapped projector family assumed in Sec.~\ref{sec:protected-response}.

\paragraph*{Two-body clustering and the Kapit--Mueller realization.}
The Laughlin state supplies the two-body clustering mechanism
\cite{laughlin1983}.  In a lowest Landau level or a Chern band, the bosonic
$V_0$ pseudopotential may be factorized as
\begin{equation}
 H_{\mathrm L}=V_0\sum_R A_R^\dagger A_R,
 \qquad
 A_R|\Psi_0\rangle=0,
 \label{eq:laughlin-parent}
\end{equation}
where $A_R$ annihilates a pair in the forbidden relative-angular-momentum
channel and $R$ labels its center-of-mass translate.  Quasihole zero modes are
common null vectors of all $A_R$; their admissible-root counting is fixed by
the $(1,2)$ clustering rule \cite{bernevighaldane2008}, and the zero-mode
property admits a purely second-quantized formulation \cite{chenseidel2015}.

Our lattice realization uses $N$ bosons on a square torus with $n_\phi$
orbitals in the exactly flat lowest Kapit--Mueller Chern band
\cite{kapitmueller2010}.  The Hilbert space is the fixed-particle-number
bosonic Fock space after projection to that band.  The boundary conditions are
magnetic-periodic; boundary twists are introduced only when the closed twist
torus is explicitly studied.  Projecting the onsite repulsion gives the
positive two-body parent in Eq.~\eqref{eq:laughlin-parent}.  The protected rank
is the generated quasihole count in Table~\ref{tab:model-evidence}, and a case
is retained only when its first positive many-body eigenvalue gives a positive
gap over the entire registered parameter set.

We use two logically distinct tangent classes.  The first is
\emph{exactness-preserving constraint/generator transport}.  A smooth
one-particle unitary $U(\boldsymbol\lambda)$, lifted to the fixed-$N$ Fock
space as $\Gamma_N[U]$, defines
$C(\boldsymbol\lambda)=C_0\Gamma_N[U(\boldsymbol\lambda)]^\dagger$ and
$H(\boldsymbol\lambda)=C(\boldsymbol\lambda)^\dagger
C(\boldsymbol\lambda)$.  Consequently $S_a=0$ along the path, the nullity is
constant, and the transported kernel remains exactly degenerate.  Statements
about finite-path parallel transport, holonomy, and protection mixing belong
to this class.

The second class comprises the \emph{v2/v3 Laughlin fixed-parent
local-potential panels}.  At the exact parent, structured Fourier and random
local potentials enter through $Q(\partial_aH)P$.  They define the linear
response of the isolated rank-$D$ cluster at the exact parent, but generally
do not preserve exact degeneracy away from the base point.  Curvature spectra
and four-channel statistics obtained from these panels are therefore
base-point response statements, not claims about a finite exact path.  An
intrafiber intervention is used only as the fixed-projector spectral control.
Together, the two classes make the Laughlin parent a spatially local,
large-rank comparison between exact transport and stochastic projector
response protected at the base point by two-body clustering.

A term with a nonscalar $P\delta HP$ breaks Laughlin exactness.  Concrete
examples are residual band dispersion, a generic projected potential
$\sum_{\bm x}\mu_{\bm x}n_{\bm x}$, or an interaction component outside the
common $V_0$ kernel.  Such terms need not destroy the surrounding fractional
Hall phase, but at finite size they remove the exact zero-energy equality that
is required here.  Truncating the ideal Kapit--Mueller hopping can likewise
replace the exact flat band by a narrow band; this distinction concerns exact
degeneracy, not phase stability.

\paragraph*{Three-body clustering and Moore--Read zero modes.}
The bosonic Moore--Read parent changes the order of the local constraint
\cite{mooreread1991}.  We work in the fixed-$N$ bosonic Fock space of
$N_\phi$ continuum lowest-Landau-level orbitals on a square torus.  On a
complete magnetic-translation orbit of coherent states, the factorized parent
is
\begin{equation}
 H_{\mathrm{MR}}=\sum_x C_x^\dagger C_x,
 \qquad
 C_x=\frac{1}{\sqrt{3!}}\left(\sum_j\phi_{xj}b_j\right)^3.
 \label{eq:moore-read-parent}
\end{equation}
Every zero mode vanishes when three particles coincide.  The fixed-four-
quasihole sequence uses $N_\phi=N+2$ with periodic boundary conditions.  Its
protected rank is obtained from the cyclic $(2,2)$ admissible count and
checked against the numerical nullity; the corresponding generated ranks and
positive gaps appear in Table~\ref{tab:model-evidence}.  The factorized
second-quantized construction is also the basis for the sparse numerical
representation \cite{zhang2023parent}.

The tangent class consists of normalized guiding-center generators organized
into local panels.  Each Hermitian one-body generator $G_a$ differentiates an
invertible transport of the three-body constraint map.  Write
$\Gamma_a=\Gamma_N(G_a)$.  On the protected fiber, the resulting tangent and
response obey
\begin{equation}
 (\partial_aH)P=-H\Gamma_aP,
 \qquad
 X_a=Q\Gamma_aP,
 \qquad S_a=0.
 \label{eq:moore-read-transport}
\end{equation}
Thus the v9 Moore--Read guiding-center panels are exactness-preserving
generator transports, even though their response is evaluated analytically
at the registered parent.  Their curvature and four-channel statistics
describe tangents to an exactly degenerate parent family, rather than generic
splitting probes.  This model's role in the comparison is to change the
clustering order and the zero-mode structure without introducing a
supercharge or a stabilizer algebra.  A generic two-body pseudopotential
$\delta H_2=\sum_R u_R A_R^\dagger A_R$, or any other term with nonscalar
$P\delta H_2P$, breaks Moore--Read exactness.  Landau-level mixing and an
incomplete three-body constraint orbit have the same effect at the level of
the ideal parent problem.

\subsection{Cohomological kernels}

\paragraph*{Random-supercharge cohomology.}
The random-supercharge model is zero dimensional: its Hilbert space is the
fermionic Fock space of $N$ complex modes, resolved into fixed-charge sectors.
There is no spatial boundary condition.  For a generic complex three-form
$C_{ijk}$, define
\begin{equation}
 \begin{aligned}
 Q(C)&=\sum_{i<j<k}C_{ijk}\psi_i\psi_j\psi_k,
 &Q^2&=0,\\
 H_{\mathrm{SUSY}}&=\{Q,Q^\dagger\}.
 \end{aligned}
 \label{eq:random-supercharge}
\end{equation}
This is the complex $\mathcal N=2$ SYK construction
\cite{fu2017susy}.  Positivity implies that the protected fiber in a charge
sector is
$\ker Q\cap\ker Q^\dagger$, equivalently the space of harmonic
representatives of $Q$-cohomology.  Its protected rank is calculated from the
charge-resolved cohomology and verified by exact nullity in every checkpoint.
Because the opened calculation contains two charge sectors and many coupling
realizations rather than one common $D$ sequence, the v13 cross-model table
leaves its single ``sampled $D$'' cell as not tested instead of selecting a
representative rank after the fact.  The positive gap is the smallest nonzero
Hodge-Laplacian eigenvalue in the same charge sector and is audited separately
for every realization.

The tangent class consists of sparse and isotropic variations $\delta C$ of
the cubic supercharge.  Differentiating $H$ separates the response into the
orthogonal exact and coexact images of $Q$ and $Q^\dagger$.  The role in the
comparison is a dense, nonlocal cohomological mechanism with an ensemble of
independent coupling realizations.  Deforming $Q$ within the nilpotent cubic
family preserves nonnegative energy and normally preserves the generic
cohomology rank, although the rank may jump at singular coupling loci.  By
contrast, adding a generic charge-preserving interaction $W$ that is not part
of an anticommutator and has nonscalar $PWP$ breaks random-supercharge
exactness by splitting the harmonic multiplet.

\paragraph*{Independence-complex cohomology on a lattice.}
For a graph $G$, let
$d_i^\dagger=c_i^\dagger\prod_{j\sim i}(1-n_j)$ create a fermion only when its
neighboring vertices are empty.  The allowed Hilbert space at charge $f$ is
therefore the vector space spanned by size-$f$ independent sets.  With complex
site couplings $z_i$,
\begin{equation}
 Q=\sum_i z_i d_i^\dagger,
 \qquad Q^2=0,
 \qquad H_{\mathrm{SUSY}}=\{Q,Q^\dagger\}.
 \label{eq:lattice-supercharge}
\end{equation}
This is the local lattice-$\mathcal N=2$ construction
\cite{fendley2003lattice}; its zero modes are computed by
independence-complex cohomology \cite{huijse2010cohomology}.

The controlled family used here is $G=\bigsqcup_{s=1}^{m}C_6$, a disjoint
union of six-site rings, in charge sector $f=2m$.  Each component has periodic
boundary conditions; different rings have no connecting edge.  The protected
rank is $2^m$ and is checked directly against the zero-eigenvalue count, while
the first positive eigenvalue supplies the gap gate.  Local tangent panels
vary the site couplings after projecting out the component-wise rescaling
directions; isotropic panels provide a matched control.  Both panels retain
the exact/coexact Hodge decomposition.  The role in the comparison is an
exactly auditable local cohomology sequence, not a connected-lattice
thermodynamic limit.

An unconstrained hopping term
$\delta H_t=\sum_{\langle ij\rangle}t_{ij}c_i^\dagger c_j+\mathrm{H.c.}$ or a
nonuniform onsite potential generally cannot be written as the same Hodge
Laplacian.  When its restriction to the harmonic fiber is nonscalar it breaks
lattice-SUSY exactness.  Altering the exclusion projector so that the proposed
supercharge is no longer nilpotent removes the cohomological protection more
directly.

\subsection{Commuting stabilizers}

The X-cube model places one qubit on every edge of an $L\times L\times L$
cubic three-torus.  Its Hilbert space has $3L^3$ qubits and periodic boundary
conditions in all three directions.  With cube $X$ stabilizers $A_c$ and
vertex-plane $Z$ stabilizers $B_v^{\mu\nu}$, the Hamiltonian is
\begin{equation}
 H_{\mathrm{XC}}=-\sum_c J_cA_c
 -\sum_{v,\mu\nu}K_v^{\mu\nu}B_v^{\mu\nu},
 \qquad J_c,K_v^{\mu\nu}>0.
 \label{eq:xcube-parent}
\end{equation}
All stabilizers commute \cite{vijay2016xcube}.  Their dependent relations on
the torus leave $6L-3$ logical qubits, hence protected rank $2^{6L-3}$.
Positive coefficients retain the same simultaneous $+1$ eigenspace and a
positive gap.  Coefficient reweighting is therefore a fixed-projector tangent
class.  A separate local unitary conjugation transports the projector while
preserving the entire spectrum; it supplies the scalar-curvature control.
The role in the comparison is an exact structured counterexample rather than
a stochastic ensemble.

Setting a stabilizer coefficient through zero closes the corresponding gap
and changes the constraint count.  Adding a logical membrane operator
$-h\overline Z$ acts nontrivially within the ground space and breaks X-cube
exactness immediately; a generic local Pauli field produces the same finite-
size splitting only at the perturbative order needed to implement a logical
operator.  This makes the distinction between a protected phase and exact
finite-size degeneracy explicit.

\subsection{Quasi-degenerate checkerboard replication}

The checkerboard calculation uses bosons in the projected lower band of a
short-range, two-sublattice $C=1$ hopping model on a periodic torus.  Boundary
twists enter the physical hopping matrix as seam phases before the band
projection.  The hopping amplitudes realize a nearly flat isolated band
\cite{sun:2011:flat}, and interactions produce an FCI-like low-energy
multiplet \cite{regnault:2011:fci}.  Bond-magnitude, bond-phase, and site-
potential disorder define the tangent path, and the onsite contact interaction
is projected into the deformed lower band at each parameter point.  The
selected low-energy cluster
has a nonzero internal bandwidth and remains separated from higher states by
a positive external gap.

Accordingly, checkerboard FCI is a quasi-degenerate replication of the local
geometric-statistics trend, not an exact-zero-mode mechanism.  Its protected
rank means the dimension of the tracked isolated cluster, not the nullity of a
frustration-free parent.  It tests sensitivity to short-range Chern-band
microscopics; it is not included when statements require exact spectral
silence.

% AUTO-GENERATED FILE. DO NOT EDIT.
% Paper I model table (v13)
% registry-payload-sha256: 95c4b13c5c9ca3c42eda11fd7bef8f0bf6c2e0f86d8a7f1233cd400f1ae0e5f1
% source centered_complete_covariance_v13: ea75047473001032c1875a02944329d3db054e788945f7882d88154a7a1ee3fc
% source continuum_lll_replication_inference_v6: 21ac615ff44ddc466a1feb5c3d6f411bde9c3d34f06fb9dc570b6e77766193b6
% source cross_complete_covariance_v12: 9f5357061d9281a25483c4315da3a09c77762091e1c6754c59a8880ec16fe0e6
% source cross_mechanism_geometric_eth_v12: 18933c332747e9a134536dc06fceed5751e5575595e315e50c2f0d93e6005b86
% source lattice_susy_pilot_v10: 1a5705686a2f37f5b06216d43d8539cf15c9d6bcb92feba74e3030661d5886e8
% source matrix_element_geometric_eth_v3: bbacad6f795559f417c1e39c034a6fe51669fd509a983099c10965fb68ee34a3
% source moore_read_pilot_v9: 3024f210ca9995f57e21c9a84750a00db274b19eb63b7ee5227add45c451071e
% source spectral_silence_v2: 5a2da65d83ded56a0fbebb23f69d5c26bd41382ebaa5f32a333fb239bffd5778
% source susy_hodge_v7_N14_inference: 177643e07fc6cf210362fc1077070bd1f0ba316b6805042a626de3f96c55a627
% source topological_holonomy_v3: ff15e7cd2f947b43da5a5973bb82c50319b164877883d187c7c46b2a073823c6
% source xcube_geometric_control_v11: c75f8a27b263d71aaf1535ed0ab54f248c67335836c3543e1b95642d1ee3cc84
% source protected_scaling_inference_v4: unavailable (not tracked in delivery branch)

\begin{table*}[t]
\caption{Mechanisms that protect exact degeneracy and the evidence currently available for geometric statistics. The centered three-pairing statistic is denoted \(R_4^{\mathrm{full}}\). Its primary gates use exact registered finite-ensemble means, a Student-\(t_{11}\) finite-sample reference, and Bonferroni control over five cases. A missing calculation is written as \(\text{not tested}\), not as a null result.}
\label{tab:model-evidence}
\begin{ruledtabular}
\begin{tabular}{@{}p{0.10\textwidth}p{0.18\textwidth}p{0.13\textwidth}p{0.09\textwidth}p{0.14\textwidth}p{0.10\textwidth}p{0.08\textwidth}@{}}
Model & Degeneracy mechanism & Geometric role & \(\substack{\text{Sampled}\\D}\) & \(\substack{R_4^{\mathrm{full}}\\\text{primary split}}\) & \(\substack{\text{Independent}\\\text{ensemble}}\) & Run state \\
\hline
% canonical column: \(R_4^{\mathrm{full}}\), primary split
% canonical Moore--Read status: N=4: not passed; N=6: passed (finite-rank)
Laughlin parents & two-body FQH parent & finite-size stochastic geometry & 16, 25, 36 & \(\text{not tested}\) & \(\text{not tested}\) & \(\text{complete}\) \\ % MODELROW:laughlin; centered-full-r4=not-tested; independent-ensemble=not-tested; production=complete
$\mathcal{N}=2$ SYK & random supercharge cohomology & finite-size stochastic geometry & \(\text{not tested}\) & \(\text{not tested}\) & \(\text{not tested}\) & \(\text{complete}\) \\ % MODELROW:susy_syk; centered-full-r4=not-tested; independent-ensemble=not-tested; production=complete
Moore--Read parent & three-body clustered FQH parent & finite-size stochastic geometry & 42, 120 & \(\substack{N=4:\\ \text{not passed}\\N=6:\\ \text{passed}\\ \text{finite rank}}\) & \(\text{not tested}\) & \(\text{pilot only}\) \\ % MODELROW:moore_read; centered-full-r4=N4-not-passed-N6-passed-finite-rank; independent-ensemble=not-tested; production=pilot-only
Lattice $\mathcal{N}=2$ SUSY & independence complex cohomology & finite-size structured geometry & 2, 4, 8 & \(\substack{m=1,2,3:\\\text{not passed}}\) & \(\text{not tested}\) & \(\text{pilot only}\) \\ % MODELROW:lattice_susy; centered-full-r4=m1-m2-m3-not-passed; independent-ensemble=not-tested; production=pilot-only
X-cube control & commuting projector subsystem code & exact structured control & \(\text{not tested}\) & \(\substack{\text{not}\\\text{applicable}}\) & \(\substack{\text{not}\\\text{applicable}}\) & \(\text{complete}\) \\ % MODELROW:xcube; centered-full-r4=not-applicable; independent-ensemble=not-applicable; production=complete
\end{tabular}
\end{ruledtabular}
\end{table*}

